\section{Related Work}

\paragraph{Agent benchmarks across environments.}
Modern agent evaluation now spans text environments, general task suites, web
interaction, enterprise work, mobile control, software engineering, and tool
use
\cite{scienceworld,agentbench,agentboard,webarena,visualwebarena,workarena,androidworld,toolgym,taubench,deepresearchbench,mcpbench,mcpuniverse}.
General-assistant and computer-use benchmarks such as \texttt{GAIA} and
\texttt{OSWorld} broaden the picture further by mixing web use, tool use, and
open-ended computer interaction under more general assistant settings
\cite{gaia,osworld}.
Repository-grounded software tasks are especially relevant for our setting
because they require understanding fixed external codebases rather than solving
pure language problems in-context. \gittaskbench joins a growing ecosystem that
includes \texttt{SWE-bench}, \texttt{SWE-agent}, \texttt{SWE-Lancer},
\texttt{Ambig-SWE}, and long software-task evaluations
\cite{swebench,sweagent,swelancer,ambigswe,longsoftwaretasks}.
Our paper does not compete with these benchmarks. It studies how one such
benchmark should be reported once official outputs have already been produced.

\paragraph{Trajectory diagnosis and post-hoc analysis.}
Several recent works argue that final success alone is insufficient for
understanding agent behavior. \texttt{AgentRx} diagnoses agent failures from
execution trajectories, while \texttt{TRAJEVAL} decomposes code-agent traces
into more granular failure modes \cite{agentrx,trajeval}. Those works motivate a
general diagnostic mindset. Our contribution is narrower: even before manual
trajectory analysis, the benchmark's own official \processmetric/\resultmetric
split already reveals useful structure.

\paragraph{Measurement blind spots in long-context evaluation.}
Long-context work has repeatedly shown that aggregate scores can hide meaningful
failure structure. Relevant examples include middle-position forgetting,
benchmark family variation, and evaluation beyond literal matching
\cite{lostmiddle,longbench,ruler,infinitebench,nolima,loogle}. We cite this
literature not because \gittaskbench is a long-context benchmark, but because it
establishes a broader measurement lesson: when a benchmark score collapses
multiple conditions of success, reporting only the aggregate can be misleading.

\paragraph{Reporting-oriented evaluation frameworks.}
Adjacent evaluation infrastructures have also argued that benchmark reports
should preserve more structure than one scalar score. \texttt{HELM} emphasizes
standardized multi-scenario and multi-metric reporting, while
\texttt{Dynabench} argues that benchmarks must stay informative even as models
quickly saturate static test sets, and \texttt{WildBench} pushes evaluation
toward harder user-originated task distributions
\cite{helm,dynabench,wildbench}. Our paper does not propose a new evaluation
platform, but it shares the same reporting instinct: once a benchmark already
exposes separable success conditions, collapsing them into one headline number
throws away evidence the evaluator has already provided.

\paragraph{Multimodal benchmarking and answer-channel sensitivity.}
Multimodal evaluation has also highlighted that benchmark outcomes depend on
what is being measured and how answers are extracted or judged. Representative
examples include \texttt{MME}, \texttt{MMMU}, \texttt{MathVista},
\texttt{HallusionBench}, and \texttt{BLINK} \cite{mme,mmmu,mathvista,hallusionbench,blink}.
The connection to our work is again methodological. Structured outputs,
thresholded criteria, and answer-channel validity matter in multimodal
benchmarks; they also matter in repository-grounded software tasks.

\paragraph{Positioning.}
This paper is best read as a measurement note on repo-grounded agent evaluation.
It does not propose a new benchmark or a new official metric. It instead argues
that \gittaskbench's existing official metrics already support more informative
reporting than a single global pass rate.
